\documentclass[aspectratio=169]{beamer}

\usetheme{Madrid}
\usecolortheme{default}

\usepackage{listings}
\usepackage{booktabs}
\usepackage{hyperref}
\usepackage{xcolor}
\usepackage{array}
\usepackage{graphicx}
\usepackage{tcolorbox}
\usepackage{tikz}

% Code listing style
\lstset{
    basicstyle=\ttfamily\small,
    breaklines=true,
    backgroundcolor=\color{gray!10},
    keywordstyle=\color{blue},
    commentstyle=\color{green!50!black},
    stringstyle=\color{red!70!black},
    showstringspaces=false
}

\title{What's New in AI: A 2026 Recap}
\subtitle{D-Lab AI Pulse Workshop}
\author[D-Lab]{Materials: \url{https://github.com/dlab-berkeley/AI_Pulse}\\Feedback form: \url{https://forms.gle/x8KfejgCViE7t6iQA}}
\date{}

\begin{document}

% ===== SLIDE 1: Title =====
\begin{frame}[plain]
    \titlepage
\end{frame}

% =============================================================================
% OPENING
% =============================================================================

% ===== SLIDE 2: Roadmap =====
\begin{frame}{Today's Roadmap}
    \begin{enumerate}
        \item What happened in 2026
        \begin{itemize}
            \item Model Context Protocol and customizable AI
            \item Agents that act on their own
            \item AI starts doing research
            \item Big tech eats every industry
        \end{itemize}
        \vspace{0.4em}
        \item What's happening right now
        \begin{itemize}
            \item Models too powerful to release
            \item AI is proving math
            \item How is AI affecting society?
        \end{itemize}
        \vspace{0.4em}
        \item What does the future look like
        \begin{itemize}
            \item The centralization question
            \item Beyond LLMs
        \end{itemize}
    \end{enumerate}
\end{frame}

% =============================================================================
% PART 1, TOPIC 1: MCP
% =============================================================================

% ===== SLIDE: MCP intro =====
\begin{frame}{MCP: Giving Tools to Your AI}
    \begin{itemize}
        \item Initially, AI was stuck in the chat window
        \vspace{0.6em}
        \item You had to continuously copy-paste, update documents, and share information
        \vspace{0.6em}
        \item The solution: Model Context Protocol, now the industry standard
        \vspace{0.6em}
        \item ``USB-C for AI'': any LLM, any tool; reads files, runs queries, takes actions
        \vspace{0.6em}
        \item Setup: thousands of ready-made MCP servers; custom ones built in a few hundred lines with the SDK
        \vspace{0.6em}
        \item Introduced by Anthropic (late 2024); now broadly supported across Claude, Cursor, VS Code, and others
    \end{itemize}
\end{frame}

% ===== SLIDE: What you can do with MCP =====
\begin{frame}{What You Can Do With MCP}
    \begin{columns}[T]
        \begin{column}{0.48\textwidth}
            Power your AI (read)
            \vspace{0.5em}
            \begin{itemize}
                \item Zotero (citations, PDFs)
                \vspace{0.4em}
                \item Google Drive / OneDrive / Box
                \vspace{0.4em}
                \item Notion / Obsidian
                \vspace{0.4em}
                \item GitHub
                \vspace{0.4em}
                \item PostgreSQL / Snowflake
                \vspace{0.4em}
                \item Local files and PDFs
            \end{itemize}
        \end{column}
        \begin{column}{0.48\textwidth}
            Let your AI act (write)
            \vspace{0.5em}
            \begin{itemize}
                \item Slack: send and search messages
                \vspace{0.4em}
                \item Asana / Linear / Jira: file tickets
                \vspace{0.4em}
                \item Gmail / Calendar: draft, schedule
                \vspace{0.4em}
                \item Figma: produce design drafts
                \vspace{0.4em}
                \item Hex and spreadsheets: analyze
                \vspace{0.4em}
                \item Browser (Playwright): navigate
            \end{itemize}
        \end{column}
    \end{columns}
\end{frame}

% ===== SLIDE: Customize and share your AI =====
\begin{frame}{Customize and Share Your AI}
    \begin{itemize}
        \item Three platforms, one idea: customize the AI, share with a link
        \vspace{0.4em}
        \item Behavior packaging: Skills (Claude), Gems (Gemini), Custom GPTs (ChatGPT)
        \vspace{0.4em}
        \item Interactive apps: Artifacts (Claude), Canvas (ChatGPT and Gemini)
        \vspace{0.4em}
        \item With Skills: format a manuscript, build slides from a dataset, run code reviews
        \vspace{0.4em}
        \item With Artifacts: build a live dashboard, simulate a physics setup, visualize math, prototype a teaching tool
    \end{itemize}
\end{frame}

% =============================================================================
% PART 1, TOPIC 2: AGENTIC AI
% =============================================================================

% ===== SLIDE: From Chats to Tasks =====
\begin{frame}{From Chats to Tasks}
    \begin{itemize}
        \item Chats: you ask, it answers
        \vspace{0.6em}
        \item Tasks: you describe an outcome; the AI plans, acts, iterates
        \vspace{0.6em}
        \item Reads your files, browses the web, runs code, edits documents
        \vspace{0.6em}
        \item Works for minutes to hours; step away, come back to results
        \vspace{0.6em}
        \item Amplifies effectiveness, but also amplifies risks
    \end{itemize}
\end{frame}

% ===== SLIDE: Claude Code and Gemini CLI =====
\begin{frame}{Claude Code and Gemini CLI}
    \begin{itemize}
        \item AI agents that live in your terminal
        \vspace{0.3em}
        \item Read and edit your codebase directly
        \vspace{0.3em}
        \item Run commands, scripts, and tests
        \vspace{0.3em}
        \item Know your file structure and Git history
        \vspace{0.3em}
        \item Spawn subagents to work in parallel
        \vspace{0.3em}
        \item Claude Code: steerable from your phone (Remote Control)
    \end{itemize}
\end{frame}

% ===== SLIDE: Cowork =====
\begin{frame}{Cowork}
    \begin{itemize}
        \item AI that runs on your desktop
        \vspace{0.3em}
        \item Reads and creates files (PDFs, Word, Excel, slides)
        \vspace{0.3em}
        \item Organizes folders, builds tables, drafts documents
        \vspace{0.3em}
        \item Connects to apps via MCP (Slack, Notion, Drive)
        \vspace{0.3em}
        \item Computer Use: controls any app on your screen (macOS research preview)
        \vspace{0.3em}
        \item Take a task, walk away, come back to results
        \vspace{0.3em}
        \item Schedule recurring jobs (daily arXiv, weekly reports)
    \end{itemize}
\end{frame}

% ===== SLIDE: Browser agents =====
\begin{frame}{Claude in Chrome (and Project Mariner, ChatGPT Agent)}
    \begin{itemize}
        \item AI that drives your browser
        \vspace{0.4em}
        \item Sees pages, clicks, fills forms, navigates tabs
        \vspace{0.4em}
        \item Library portals, journal sites, IRB forms, grant submissions
        \vspace{0.4em}
        \item Faster than full screen control for web-only tasks
        \vspace{0.4em}
        \item Equivalents: Project Mariner (Google) and ChatGPT Agent (OpenAI)
    \end{itemize}
\end{frame}

% ===== SLIDE: OpenClaw =====
\begin{frame}{OpenClaw}
    \begin{itemize}
        \item Open-source agent in your messaging apps (Signal, Telegram, WhatsApp, Discord)
        \vspace{0.4em}
        \item 100+ skills; model-agnostic; runs locally on your machine
        \vspace{0.4em}
        \item Common uses: morning briefings, email triage, coding from your phone
        \vspace{0.4em}
        \item Power workflows: parallel research, smart home via chat, self-healing servers
        \vspace{0.4em}
        \item Viral arc: 250K stars in 60 days; creator joined OpenAI Feb 2026
        \vspace{0.4em}
        \item Caveat: Cisco found 12\% of community skills compromised
        \vspace{0.4em}
        \item Highlights growing security risks in open agent ecosystems
    \end{itemize}
\end{frame}

% =============================================================================
% PART 1, TOPIC 3: AUTOMATED RESEARCH
% =============================================================================

% ===== SLIDE: Automated research transition =====
\begin{frame}{Automated Research: Agents Doing Science}
    \begin{itemize}
        \item Agents can read your files and act on your tools
        \vspace{0.3em}
        \item The next question: can they do science?
        \vspace{0.3em}
        \item Research is a loop: hypothesize, experiment, analyze, write, review
        \vspace{0.3em}
        \item In 2026, every step has a working AI tool
        \vspace{0.3em}
        \item Important caveat: AI-generated papers and experiments still require rigorous human verification
        \vspace{0.3em}
        \item Typical cost: tens to hundreds of dollars per full research run
    \end{itemize}
\end{frame}

% ===== SLIDE: Karpathy AutoResearch =====
\begin{frame}{Karpathy's AutoResearch}
    \begin{itemize}
        \item Andrej Karpathy, released March 2026
        \vspace{0.3em}
        \item $\sim$630 lines of Python; tens of thousands of GitHub stars in days
        \vspace{0.3em}
        \item AI agent runs a small ML training setup
        \vspace{0.3em}
        \item Modifies code, trains 5 minutes, scores result, keeps or discards, repeats
        \vspace{0.3em}
        \item 700 experiments in 2 days; 20 optimizations discovered
        \vspace{0.3em}
        \item Karpathy's term: ``the self-improvement loopy era''
    \end{itemize}
    \vspace{0.4em}
    \footnotesize \url{github.com/karpathy/autoresearch}
\end{frame}

% ===== SLIDE: Sakana AI Scientist =====
\begin{frame}{Sakana AI Scientist}
    \begin{itemize}
        \item End-to-end agent: idea, code, experiments, paper, peer review
        \vspace{0.3em}
        \item First AI-generated paper to pass double-blind peer review at the ICLR 2025 ICBINB workshop; later withdrawn as part of the experiment
        \vspace{0.3em}
        \item Paper: ``Compositional Regularization: Unexpected Obstacles in Enhancing Neural Network Generalization''
        \vspace{0.3em}
        \item Sakana team's human-authored Nature paper (March 2026) describes the system and documents a clear scaling law: better base models produce dramatically better AI-generated papers
    \end{itemize}
    \vspace{0.4em}
    \footnotesize \url{github.com/SakanaAI/AI-Scientist-v2} \quad \url{sakana.ai/ai-scientist-first-publication}
\end{frame}

% ===== SLIDE: HKU AI-Researcher =====
\begin{frame}{HKU AI-Researcher}
    \begin{itemize}
        \item NeurIPS 2025 Spotlight; University of Hong Kong
        \vspace{0.3em}
        \item Reads relevant literature
        \vspace{0.3em}
        \item Generates novel research ideas
        \vspace{0.3em}
        \item Codes and runs experiments
        \vspace{0.3em}
        \item Analyzes results and writes the manuscript
        \vspace{0.3em}
        \item Live demo at novix.science
    \end{itemize}
    \vspace{0.4em}
    \footnotesize \url{github.com/HKUDS/AI-Researcher}
\end{frame}

% ===== SLIDE: refine.ink =====
\begin{frame}{refine.ink}
    \begin{itemize}
        \item AI as an obsessive peer reviewer
        \vspace{0.3em}
        \item Founded 2025 by Benjamin Golub (Northwestern) and Yann Calvo Lopez (CNRS)
        \vspace{0.3em}
        \item Reads finished drafts; flags errors, inconsistencies, logical flaws
        \vspace{0.3em}
        \item 30-minute turnaround using hundreds of parallel LLM calls
        \vspace{0.3em}
        \item Used at Harvard, Stanford, Yale, MIT, Caltech
        \vspace{0.3em}
        \item Closed source; first review free
    \end{itemize}
\end{frame}

% =============================================================================
% PART 1, TOPIC 4: ANTHROPIC DISRUPTING INDUSTRIES
% =============================================================================

% ===== SLIDE: Boris Cherny tweet =====
\begin{frame}{How Anthropic Uses Its Own Tools}
    \begin{tcolorbox}[colback=blue!5,colframe=blue!40,boxrule=0.8pt,arc=4pt]
    \textit{``I'm Boris and I created Claude Code. Lots of people have asked how I use Claude Code...''}

    \vspace{0.4em}
    \textit{``I run 5 Claudes in parallel in my terminal. I number my tabs 1-5, and use system notifications to know when a Claude needs input.''}

    \vspace{0.4em}
    \textit{``I also run 5-10 Claudes on claude.ai/code, in parallel with my local Claudes. I also start a few sessions from my phone every morning and throughout the day.''}
    \end{tcolorbox}

    \vspace{0.6em}
    Boris Cherny (creator of Claude Code), January 2, 2026
\end{frame}

% =============================================================================
% PART 2, TOPIC 1: AI BEHIND CLOSED DOORS
% =============================================================================

% ===== SLIDE: Anthropic Mythos =====
\begin{frame}{Anthropic Mythos}
    \begin{itemize}
        \item Anthropic's newest frontier model: a step up in general capability
        \vspace{0.4em}
        \item Particularly strong at finding software flaws (downstream of better coding and reasoning)
        \vspace{0.4em}
        \item So powerful Anthropic delayed public release over dual-use concerns
        \vspace{0.4em}
        \item Project Glasswing: access for security companies first, so defenders can harden systems
        \vspace{0.4em}
        \item Found vulnerabilities in FFmpeg, Firefox, and FreeBSD
    \end{itemize}
\end{frame}

% ===== SLIDE: Supply chain twist =====
\begin{frame}{The Supply Chain Twist}
    \begin{itemize}
        \item Mar 3, 2026: DoD designates Anthropic a supply-chain risk; Pentagon blocked from Mythos
        \vspace{0.3em}
        \item The NSA is using Mythos anyway
        \vspace{0.3em}
        \item Apr 16: White House moves to give other US agencies access (Bloomberg)
        \vspace{0.3em}
        \item Banks pursuing access: Goldman, Citi, BofA, Morgan Stanley
        \vspace{0.3em}
        \item China's 360 Digital Security has a parallel tool ($\sim$1{,}000 unknown bugs found)
        \vspace{0.3em}
        \item Google M-Trends 2026: mean time to exploit has gone negative
    \end{itemize}
\end{frame}

% =============================================================================
% PART 2, TOPIC 2: AI FOR MATH
% =============================================================================

% ===== SLIDE: AI can now reason (transition) =====
\begin{frame}{AI Can Now Reason}
    \begin{itemize}
        \item The big unlock of 2026: AI is doing real mathematical reasoning
        \vspace{0.3em}
        \item Not just generating plausible-looking proofs
        \vspace{0.3em}
        \item Actually closing open problems
        \vspace{0.3em}
        \item Math is the cleanest test: a proof either works or it doesn't
        \vspace{0.3em}
        \item Speed of discoveries is accelerating
    \end{itemize}
\end{frame}

% ===== SLIDE: Timeline =====
\begin{frame}{From One Problem to One a Day}
    \begin{itemize}
        \item Oct 2025: GPT-5 helps Sawhney + Sellke rediscover forgotten 2003 solution to Erd\H{o}s \#339 via superior literature retrieval
        \vspace{0.3em}
        \item Jan 2026: First fully autonomous Lean proof (Aristotle agent, \#728)
        \vspace{0.3em}
        \item Apr 9: OpenAI internal model closes 5 problems in a single batch
        \vspace{0.3em}
        \item Apr 24, 2026: Liam Price solves Erd\H{o}s \#1196 in a single prompt (raw output required significant expert refinement)
        \vspace{0.3em}
        \item Apr 25, 2026: Additional progress on Erd\H{o}s \#1138 using GPT-5.5
    \end{itemize}
\end{frame}

% ===== SLIDE: AI as math collaborator (merged) =====
\begin{frame}{AI as Math Collaborator}
    \begin{itemize}
        \item AI is now part of math research at every level: experts and amateurs alike
        \vspace{0.5em}
        \item Mehtaab Sawhney (Columbia) and Mark Sellke (Harvard): used GPT-5 to surface a forgotten 2003 result and crack Erd\H{o}s \#339
        \vspace{0.5em}
        \item Liam Price (23-year-old amateur): cracked Erd\H{o}s \#1196 in a single ChatGPT prompt
        \vspace{0.5em}
        \item Pattern: AI proposes ideas and partial proofs; humans formalize and verify
        \vspace{0.5em}
        \item Tao: ``The job description is changing''
    \end{itemize}
    \vspace{0.4em}
    \footnotesize \url{https://chatgpt.com/share/69dd1c83-b164-8385-bf2e-8533e9baba9c}
\end{frame}

% =============================================================================
% PART 2, TOPIC 3: HOW IS AI AFFECTING SOCIETY?
% =============================================================================

% ===== SLIDE: How is AI affecting society (transition) =====
\begin{frame}{How Is AI Affecting Society?}
    \begin{itemize}
        \item Beyond the labs, AI is reshaping society in real time
        \vspace{0.5em}
        \item The loudest debate: jobs and the labor market
        \vspace{0.5em}
        \item Quieter but big: defense entanglement, community pushback, energy demand
    \end{itemize}
\end{frame}

% ===== SLIDE: Dario's Warning =====
\begin{frame}{Dario Amodei's Warning}
    \begin{itemize}
        \item Anthropic CEO; Axios interview, May 2025
        \vspace{0.3em}
        \item Up to 50\% of entry-level white-collar jobs could disappear
        \vspace{0.3em}
        \item US unemployment to 10-20\% in 1-5 years
        \vspace{0.3em}
        \item Sectors hit hardest: tech, finance, law, consulting
        \vspace{0.3em}
        \item ``Cancer is cured, the economy grows at 10\% a year, the budget is balanced, and 20\% of people don't have jobs.''
        \vspace{0.3em}
        \item Continued to warn strongly throughout 2026, including at Davos; software engineers replaced in 6-12 months
    \end{itemize}
\end{frame}

% ===== SLIDE: Companies cutting =====
\begin{frame}{Companies Already Cutting}
    \begin{itemize}
        \item Salesforce (Benioff, Jan 2025): ``maybe we aren't going to hire anybody this year''; cited 30\%+ engineering productivity
        \vspace{0.3em}
        \item Klarna (May 2025): reversed its AI-only customer-service push after quality dropped
        \vspace{0.3em}
        \item Wall Street Q1 2026: 6 big banks cut 5{,}000 net (~15{,}000 gross) while booking \$47B profit
        \vspace{0.3em}
        \item Big Tech Q1 2026: Meta + Microsoft $\sim$20{,}000 combined; Amazon $\sim$30{,}000 since Oct 2025
        \vspace{0.3em}
        \item AI cited in 25\% of March 2026 layoff announcements (Challenger, Gray \& Christmas)
    \end{itemize}
\end{frame}

% ===== SLIDE: The Dissenters =====
\begin{frame}{The Dissenters}
    \begin{itemize}
        \item Yann LeCun (April 2026): ``Dario is wrong. He knows absolutely nothing about the effects of technological revolutions on the labor market.''
        \vspace{0.3em}
        \item LinkedIn (April 15): hiring down 20\% since 2022, but no measurable AI displacement signal in support, admin, or marketing
        \vspace{0.3em}
        \item Andrew Ng (Davos 2026): ``AI can only do 30-40\% of the work now and for the foreseeable future''
        \vspace{0.3em}
        \item Daron Acemoglu (MIT, Nobel 2024): only $\sim$20\% of tasks AI-exposed; modest productivity gains
        \vspace{0.3em}
        \item Sundar Pichai (Cloud Next 2026): augmentation-positive; ``every employee can become a builder''
    \end{itemize}
\end{frame}

% ===== SLIDE: Beyond jobs =====
\begin{frame}{Beyond Jobs: AI's Broader Footprint}
    \begin{itemize}
        \item Anthropic Institute (March 2026) houses the Societal Impacts team; OpenAI's parallel is its Global Affairs and Policy team
        \vspace{0.4em}
        \item Pentagon-Anthropic row: DoD designated Anthropic a supply-chain risk (March 2026); federal judge blocked the designation
        \vspace{0.4em}
        \item Palantir defense: Maven targeting in Iran; ImmigrationOS for ICE deportations; \$300M USDA food-supply deal (April 22)
        \vspace{0.4em}
        \item Communities pushing back: Maine's legislature passed a statewide data center moratorium (later vetoed by governor)
        \vspace{0.4em}
        \item Utah: Kevin O'Leary's 9 GW Wonder Valley campus approved, projected to consume more than twice the state's current electricity
    \end{itemize}
\end{frame}

% =============================================================================
% PART 3: WHAT DOES THE FUTURE LOOK LIKE
% =============================================================================

% =============================================================================
% PART 3, TOPIC 1: ALTMAN ON CENTRALIZATION
% =============================================================================

% ===== SLIDE: Altman on centralization =====
\begin{frame}{Sam Altman: Centralization Could Lead to Ruin}
    \begin{itemize}
        \item Feb 19, 2026 (India AI Impact Summit): ``Centralisation of this technology in one company or country could lead to ruin''
        \vspace{0.3em}
        \item Called for an IAEA-style global AI regulator
        \vspace{0.3em}
        \item Warned superintelligence could arrive ``within a couple of years''
        \vspace{0.3em}
        \item ``The Gentle Singularity'' (June 2025): ``Make superintelligence cheap, widely available, and not too concentrated''
        \vspace{0.3em}
        \item The tension: OpenAI is simultaneously building ``a gigawatt of new AI infrastructure every week'' (BlackRock Summit, March 2026)
    \end{itemize}
\end{frame}

% =============================================================================
% PART 3, TOPIC 2: WORLD MODELS AND INEFFABLE INTELLIGENCE
% =============================================================================

% ===== SLIDE: LeCun =====
\begin{frame}{Yann LeCun: LLMs Are a Dead End}
    \begin{itemize}
        \item ``LLMs are an off-ramp on the road to human-level AI''
        \vspace{0.3em}
        \item Nov 2025: left Meta after 12 years (FAIR's founding director)
        \vspace{0.3em}
        \item March 2026: founded AMI Labs (Advanced Machine Intelligence)
        \vspace{0.3em}
        \item Raised \$1.03B at \$3.5B pre-money
        \vspace{0.3em}
        \item His architecture: JEPA, predicting representations rather than tokens
        \vspace{0.3em}
        \item He's not alone
    \end{itemize}
\end{frame}

% ===== SLIDE: Fei-Fei Li and World Models =====
\begin{frame}{Fei-Fei Li and World Models}
    \begin{itemize}
        \item Founded World Labs in 2024 with Justin Johnson, Christoph Lassner, Ben Mildenhall
        \vspace{0.3em}
        \item Thesis: ``spatial intelligence''; AI must understand 3D space, physics, and persistent geometry
        \vspace{0.3em}
        \item Marble (launched Nov 2025): generates persistent, editable 3D environments from text, images, video
        \vspace{0.3em}
        \item Feb 2026: \$1B from AMD, Autodesk, NVIDIA, Fidelity; in talks at \$5B valuation
        \vspace{0.3em}
        \item The bet: language alone isn't enough; physical understanding is the unlock
    \end{itemize}
\end{frame}

% ===== SLIDE: Ineffable Intelligence =====
\begin{frame}{Ineffable Intelligence}
    \begin{itemize}
        \item Founded by David Silver (DeepMind, architect of AlphaGo and AlphaZero)
        \vspace{0.3em}
        \item April 27, 2026 (yesterday): \$1.1B seed at \$5.1B valuation
        \vspace{0.3em}
        \item Sequoia, NVIDIA, UK government; largest European seed ever
        \vspace{0.3em}
        \item Thesis: a ``superlearner'' that learns from environmental interaction with no pre-training on human data
        \vspace{0.3em}
        \item Foundation paper: Silver and Richard Sutton (2024 ACM Turing Award co-recipient), ``Welcome to the Era of Experience''
    \end{itemize}
\end{frame}

% ===== SLIDE: What they promise =====
\begin{frame}{What They Promise}
    \begin{itemize}
        \item Predict how the world changes after an action (physical reasoning)
        \vspace{0.5em}
        \item Robotics and embodied control (Gemini Robotics, NVIDIA Cosmos)
        \vspace{0.5em}
        \item Persistent 3D worlds (Marble, DeepMind's Genie 3 at 24 fps)
        \vspace{0.5em}
        \item Discovery beyond human data (AlphaProof's IMO silver via self-play)
        \vspace{0.5em}
        \item Three Turing-adjacent figures betting against pure LLM scaling
    \end{itemize}
\end{frame}

\end{document}
